\subsection{Hilbert-Schmidt 映射}

定义 9. --- 设 E 与 F 为两个希尔伯特空间。从 E 到 F 中的连续线性映射 u 称为 Hilbert-Schmidt 映射，如果 E 的正自同态 \(u^{*}u\) 的迹是有限的。用 \(\mathcal{L}^{2}(\mathrm{E},\mathrm{F})\) 表示从 E 到 F 中的所有 Hilbert-Schmidt 映射组成的集。

当 \(\mathrm{E} = \mathrm{F}\) 时，我们把 \(\mathcal{L}^2(\mathrm{E})\) 用作 \(\mathcal{L}^2(\mathrm{E};\mathrm{E})\) 的记号。

对每个 \(u \in \mathcal{L}(\mathrm{E},\mathrm{F})\)，令 \(\| u \|_2 = \operatorname{Tr}(u^* u)^{1/2}\)，于是 \(u\) 属于 \(\mathcal{L}^2(\mathrm{E};\mathrm{F})\) 当且仅当 \(\| u \|_2\) 是有限的。由迹的定义，我们得到

\[
\| u \| _ {2} ^ {2} = \sup _ {x _ {1}, \dots , x _ {m}} \sum_ {i = 1} ^ {m} \left\| u (x _ {i}) \right\| ^ {2}\tag{32}
\]

其中 \((x_{1},\ldots ,x_{m})\) 取遍 E 中的有限标准正交序列组成的集。特别地，在公式 (32) 中取 \(m = 1\)，我们有

\[
\| u \| \leqslant \| u \| _ {2} \quad (u \in \mathscr {L} (\mathrm{E}; \mathrm{F})).\tag{33}
\]

设 \((e_i)_{i\in I}\) 为 \(E\) 的标准正交基，\((f_j)_{j\in J}\) 为 \(F\) 的标准正交基。由 \(V\), p.~50 的公式 (25) 及 Parseval 等式（V, p.~22），我们有

\[
\left\| u \right\| _ {2} ^ {2} = \sum_ {i \in I} \left\| u (e _ {i}) \right\| ^ {2} = \sum_ {i, j} \left| \langle f _ {j} | u (e _ {i}) \rangle \right| ^ {2}.\tag{34}
\]

由于 \(\left|\langle f_j | u(e_i) \rangle\right| = \left|\langle e_i | u^*(f_j) \rangle\right|\)，公式 (34) 蕴含

\[
\| u ^ {*} \| _ {2} = \| u \|;\tag{35}
\]

因此，Hilbert-Schmidt 映射的伴随映射仍是 Hilbert-Schmidt 映射。设 \(E_{1}\)、\(F_{1}\) 为希尔伯特空间，\(v: E_{1} \to E\)、\(w: F \to F_{1}\) 为连续线性映射。由 (32)，我们立即导出

\[
\| w u \| _ {2} \leqslant \| w \|. \| u \| _ {2}.\tag{36}
\]

由 (35)、(36) 及关系式 \(uv = (v^{*}u^{*})^{*}\)，我们得到

\[
\left\| u v \right\| _ {2} \leqslant \left\| u \right\| _ {2} \left\| v \right\|.\tag{37}
\]

特别地，若 \(u\) 属于 \(\mathcal{L}^2(\mathrm{E},\mathrm{F})\)，则 wuv 属于 \(\mathcal{L}^2(\mathrm{E}_1,\mathrm{F}_1)\)。

定理 1. --- 设 E 与 F 为两个希尔伯特空间。

\begin{enumerate}
\def\labelenumi{(\roman{enumi})}
\item
  集 \(\mathcal{L}^2 (\mathrm{E},\mathrm{F})\) 是 \(\mathcal{L}(\mathrm{E};\mathrm{F})\) 的向量子空间，且 \(u\mapsto \| u\| _2\) 是 \(\mathcal{L}^2 (\mathrm{E};\mathrm{F})\) 上的一个希尔伯特范数（V, p.~6）。
\item
  同构 \(\theta\)（从 \(\mathbf{F} \otimes \mathbf{E}'\) 到 \(\mathcal{L}_f(\mathbf{E};\mathbf{F})\) 上，由 \(\theta(y \otimes x^*) = yx^*\) 刻画）可延拓为一个同构 \(\hat{\theta}\)（从 \(\mathbf{F} \hat{\otimes}_2 \mathbf{E}'\) 到 \(\mathcal{L}^2(\mathbf{E};\mathbf{F})\) 上）。特别地，\(\mathcal{L}_f(\mathbf{E};\mathbf{F})\) 在 \(\mathcal{L}^2(\mathbf{E};\mathbf{F})\) 中稠密。
\end{enumerate}

设 \((e_i)_{i\in \mathbf{I}}(\mathrm{resp.}(f_j)_{j\in \mathbf{J})}\) 为 E（相应地 F）的标准正交基。对每个 \(u\in \mathcal{L}(\mathbf{E};\mathbf{F})\)，令 \(\Lambda (u)\) 为 \(u\) 关于 E 与 F 的所选标准正交基的矩阵（V, p.~22）。用 \(\| \mathbf{a}\| _2\) 表示元素 \(a\)（属于希尔伯特空间 \(\ell^2 (\mathbf{J}\times \mathbf{I})\)）的范数。由公式 (34)，\(\Lambda\) 是一个从 \(\mathcal{L}^2 (\mathbf{E};\mathbf{F})\) 到 \(\ell^2 (\mathbf{J}\times \mathbf{I})\) 中的映射，使得 \(\| \Lambda (u)\| _2 = \| u\|\)；显然 \(\Lambda\) 是单射。为证明 (i)，只需证明 \(\Lambda\) 是满射。设 \(\mathbf{a} = (a_{ji})\) 为 \(\ell^2 (\mathbf{J}\times \mathbf{I})\) 的一个元素；由 Cauchy-Schwarz 不等式，我们有

\[
\left| \sum_ {j, i} \overline {{{{\eta}}}} _ {j} a _ {j i} \xi_ {i} \right| ^ {2} \leqslant \sum_ {j, i} | a _ {j i} | ^ {2} \sum_ {j, i} | \overline {{{{\eta}}}} _ {j} \xi_ {i} | ^ {2} = \| \mathbf {a} \| _ {2} ^ {2} \| \boldsymbol {\xi} \| ^ {2} \| \boldsymbol {\eta} \| ^ {2}
\]

（对每个 \(\xi = (\xi_i)\)（属于 \(\ell^2 (\mathrm{I})\)）及 \(\pmb{\eta} = (\eta_j)\)（属于 \(\ell^2 (\mathrm{J})\)））。则存在连续半双线性型 \(\Phi\)（在 \(\mathbf{F}\times \mathbf{E}\) 上），使得 \(\Phi (y,x) = \sum_{j,i}\overline{\eta}_ja_{ji}\xi_i\)（对 E 中的 \(x = \sum_{i}\xi_{i}e_{i}\) 及 F 中的 \(y = \sum_{j}\eta_{j}f_{j}\)）。设 \(u\in \mathcal{L}(\mathrm{E};\mathrm{F})\) 使得 \(\Phi (y,x) = \langle y|u(x)\rangle\)（V, p.~16, 推论 2）。我们得到

\[
a _ {j i} = \Phi (f _ {j}, e _ {i}) = \langle f _ {j} | u (e _ {i}) \rangle \quad \text { for } \quad i \in \mathrm{I}, j \in \mathrm{J}  ,
\]

因而 \(\mathbf{a} = \Lambda(u)\)。

由于 \(\Lambda\) 是从 \(\mathcal{L}^2 (\mathrm{E};\mathrm{F})\) 到 \(\ell^2 (\mathrm{J}\times \mathrm{I})\) 上的希尔伯特空间同构，且 \((f_j\otimes e_i^*)\) 是 \(\mathbf{F}\hat{\otimes}_2\mathbf{E}'\) 的标准正交基，故存在一个同构 \(\hat{\theta}\)（从 \(\mathbf{F}\hat{\otimes}_2\mathbf{E}'\) 到 \(\mathcal{L}^2 (\mathrm{E};\mathrm{F})\) 上），使得

\[
\langle f _ {j} | \hat {\theta} (t) e _ {i} \rangle = \langle f _ {j} \otimes e _ {i} ^ {*} | t \rangle
\]

（对每个 \(i \in \mathbf{I}\)、\(j \in \mathbf{I}\) 及 \(t \in \mathbf{F} \hat{\otimes}_2 \mathbf{E}'\)）。特别地，取 \(t = y \otimes x^*\)，我们得到

\[
\langle f _ {j} | \hat {\theta} (y \otimes x ^ {*}) e _ {i} \rangle = \langle f _ {j} \otimes e _ {i} ^ {*} | y \otimes x ^ {*} \rangle = \langle f _ {j} | y \rangle \langle x | e _ {i} \rangle = \langle f _ {j} | y x ^ {*} e _ {i} \rangle
\]

因而 \(\hat{\theta}(y\otimes x^{*}) = yx^{*}\)。这就证明了 (ii)。

证毕。

例. --- 1) 设 I 与 J 为两个集。由上面给出的证明，为使映射 \(u\)（从 \(\ell^2(\mathbf{I})\) 到 \(\ell^2(\mathbf{J})\) 中）是 Hilbert-Schmidt 映射，必要且充分的是：存在一个矩阵 \((a_{ji})\)（属于 \(\ell^2(\mathbf{J} \times \mathbf{I})\)），使得 \(u(\xi)_j = \sum_{i \in \mathbf{I}} a_{ji} \xi_i\) 对所有 \(\xi = (\xi_i)\)（属于 \(\ell^2(\mathbf{I})\)）成立。

* 2) 设 X 与 Y 为两个 Hausdorff 拓扑空间，分别赋有正测度 \(\mu\) 与 \(\nu\)。我们可以证明，从 \(\mathcal{L}^2(\mathrm{X})\) 到 \(\mathcal{L}^2(\mathrm{Y})\) 中的 Hilbert-Schmidt 映射与 \(\mathrm{Y} \times \mathrm{X}\) 上的平方可积函数的类一一对应；与函数 \(\mathrm{N} \in \mathcal{L}^2(\mathrm{Y} \times \mathrm{X}, \nu \otimes \mu)\) 的类对应的是由下式给出的映射 \(u_{\mathrm{N}}\)

\[
(u _ {\mathrm{N}} f) (y) = \int_ {\mathrm{X}} \mathrm{N} (y, x) f (x) d \mu (x)\tag{38}
\]

（对 v-几乎所有的 \(y \in \mathbf{Y}\) 及 \(f \in \mathcal{L}^2(\mathbf{X}, \mu)\)）。我们有

\[
\| u _ {\mathrm{N}} \| _ {2} ^ {2} = \int_ {\mathrm{X}} \int_ {\mathrm{Y}} | \mathrm{N} (y, x) | ^ {2} d \mu (x) d v (y). *\tag{39}
\]

注. --- 1) 设 \(\mathbf{K} = \mathbf{C}\)。设 \(u\) 与 \(v\) 属于 \(\mathcal{L}^2(\mathrm{E};\mathrm{F})\)。我们有关系式 \(4u^{*}v = \sum_{\varepsilon^{4} = 1}\overline{\varepsilon} (u + \varepsilon v)^{*}(u + \varepsilon v)\)，因而 \(u^{*}v\) 属于 \(\mathcal{L}^1(\mathrm{E})\)。希尔伯特空间 \(\mathcal{L}^2(\mathrm{E};\mathrm{F})\) 中的内积由下式给出

\[
\langle u | v \rangle = \operatorname{Tr} (u ^ {*} v)\tag{40}
\]

因为这一公式在 \(\mathcal{L}^{2}(\mathrm{E};\mathrm{F})\) 上定义了一个埃尔米特型，并且我们得到 \(\langle u|u\rangle=\|u\|_{2}^{2}\)。若 \(u\in\mathcal{L}^{2}(\mathrm{E};\mathrm{F})\) 且 \(v\in\mathcal{L}^{2}(\mathrm{F};\mathrm{E})\)，则由前面的结果，vu 属于 \(\mathcal{L}^{1}(\mathrm{E})\)，uv 属于 \(\mathcal{L}^{1}(\mathrm{F})\)；此外，我们有

\[
\operatorname{Tr} (u v) = \operatorname{Tr} (v u).\tag{41}
\]

由线性性与连续性，只需在 \(u = y_{1}x_{1}^{*}\) 且 \(v = x_{2}y_{2}^{*}\)（其中 \(x_{1}, x_{2}\) 属于 E，\(y_{1}, y_{2}\) 属于 F）时验证这一公式；但这时 uv 是映射 \(y \mapsto y_{1} \langle x_{1}|x_{2} \rangle \langle y_{2}|y \rangle\)，而 vu 是映射 \(x \mapsto x_{2} \langle y_{2}|y_{1} \rangle \langle x_{1}|x \rangle\)，于是 (41) 由 V, p.~48 的公式 (22) 得出。

因此，若 \(u_{1}, u_{2}\) 是 \(\mathcal{L}^2(\mathrm{E};\mathrm{F})\) 的两个元素，则我们在希尔伯特空间 \(\mathcal{L}^2(\mathrm{F};\mathrm{E})\) 中有

\[
\langle u _ {1} ^ {*} | u _ {2} ^ {*} \rangle = \operatorname{Tr} (u _ {1} u _ {2} ^ {*}) = \operatorname{Tr} (u _ {2} ^ {*} u _ {1}) = \langle u _ {2} | u _ {1} \rangle = \overline {{\langle u _ {1} | u _ {2} \rangle}};\tag{42}
\]

换言之，\(u \mapsto u^*\) 是从希尔伯特空间 \(\mathcal{L}^2(\mathrm{E};\mathrm{F})\) 到希尔伯特空间 \(\mathcal{L}^2(\mathrm{F};\mathrm{E})\) 的共轭空间（V, p.~6）上的同构。若把这个共轭空间与 \(\mathcal{L}^2 (\mathrm{F};\mathrm{E})\) 的对偶（V, p.~15）等同起来，我们就看到 \(\mathcal{L}^2 (\mathrm{E};\mathrm{F})\) 可与 \(\mathcal{L}^2 (\mathrm{F};\mathrm{E})\) 的对偶等同，典范双线性型 \((v,u)\mapsto \langle v,u\rangle\) 等同于 \((v,u)\mapsto \operatorname {Tr}(vu)\)。

\begin{enumerate}
\def\labelenumi{\arabic{enumi})}
\setcounter{enumi}{1}
\tightlist
\item
  设 K = R。我们留给读者去验证公式 (40) 与 (41) 仍然有效，并证明 \(\mathcal{L}^{2}(E; F)\) 可与 \(\mathcal{L}^{2}(F; E)\) 的对偶借助双线性型 \((u, v) \mapsto \operatorname{Tr}(uv)\) 等同起来。
\end{enumerate}

